%%
%%
%%
\documentclass[a4paper,12pt]{article}

\usepackage{a4wide}
\usepackage{wasysym}
\usepackage{times}
\usepackage{framed}
\usepackage{graphics}
\usepackage{graphicx}

\title{Esercitazione di Embedded Systems}
\date{3 Novembre 2025}
\author{prof. Corrado Santoro}

\begin{document}

\maketitle
\def\X{\textbf{``X''}}
\def\Y{\textbf{``Y''}}
\def\Z{\textbf{``Z''}}
\def\T{\textbf{``T''}}

%%
%%\vspace{1em}
%%
Realizzare un programma per microcontroller STM32 che simuli un sistema
di movimentazione di un asse lineare
secondo le seguenti specifiche.

Si supponga un asse lineare di $1~m$ che \`e possibile posizionare
con la risoluzione di $1~cm$.

Le prime due cifre del display devono mostrare la posizione attuale
(nell'intervallo $[0, 99]$), mentre le cifre successive devono mostrare
la posizione desiderata, ``target'', che si vuole raggiungere.

Il sistema deve operare nel seguente modo:
\begin{itemize}
\item tasto \X{}, incrementa la posizione target di $1~cm$;
\item tasto \Y{}, decrementa la posizione target di $1~cm$;
\item tasto \Z{}, avvia il movimento;
\item tasto \T{}, imposta la posizione target a $0$ e avvia il movimento.
\end{itemize}

Il movimento deve essere simulato usando un apposito timer e considerando
una velocita' uniforme di $5~cm/s$.
Durante il moto, il LED rosso deve lampeggiare (cambiare stato) ogni volta che l'asse si sposta di $1~cm$.
In nessun caso l'asse deve superare i limiti dell'intervallo $[0,99]$.

%%
\end{document}
%%
